\chapter{Conception fonctionnelle et technique}

\textbf{Note :} Cette phase de conception est cruciale et a été complètement terminée avant le démarrage du développement. Pour \textbf{My Smart Budget}, j'ai construit une conception fonctionnelle et technique détaillée, validée itérativement, qui sert de référence unique au développement, aux tests et à la soutenance.

\section{Méthodologie de conception}

La conception d'\textbf{My Smart Budget} s'est déroulée en plusieurs étapes structurées :

\begin{enumerate}
    \item \textbf{Cadrage fonctionnel :} Identification des besoins critiques à partir des personae.
    \item \textbf{Modélisation des Use Cases :} Liste et priorisation des cas d'usage (MoSCoW).
    \item \textbf{Diagrammes UML :} Cas d'utilisation, Classes et Séquences pour les flux critiques.
    \item \textbf{Conception UI/UX :} Zoning, Wireframes et Maquettes Figma haute fidélité.
    \item \textbf{Modélisation des données :} MCD $\rightarrow$ MLD $\rightarrow$ MPD (PostgreSQL).
    \item \textbf{Architecture 3-tiers :} Formalisation de la stack Next.js / Express / PostgreSQL + MongoDB.
    \item \textbf{Plan de tests :} Définition des critères d'acceptation par User Story.
    \item \textbf{Validation :} Relecture globale avant implémentation.
\end{enumerate}


% ============================================================
% USE CASES
% ============================================================

\section{Modélisation fonctionnelle (UML)}

\subsection{Acteurs et Cas d'Usage}

Les interactions du système sont définies par trois acteurs principaux :
\begin{itemize}
    \item \textbf{Utilisateur authentifié :} Suit son budget (Léa, Marc).
    \item \textbf{Administrateur :} Monitore l'application et les statistiques globales.
    \item \textbf{Système externe :} Service d'envoi d'emails et notifications.
\end{itemize}

\textbf{Principaux cas d'usage (UC) :}
\begin{itemize}
    \item \textbf{UC01 -- Compte :} Inscription, connexion, gestion profil. Dans ce diagramme, c'est le point d'entrée de tout utilisateur dans l'application.
    \item \textbf{UC02 -- Transactions :} CRUD complet, filtres (Revenus/Dépenses). Cas central : l'utilisateur ajoute une dépense et voit son budget mis à jour en temps réel.
    \item \textbf{UC03 -- Enveloppes :} Définition des plafonds, suivi consommation.
    \item \textbf{UC04 -- Dashboard :} Visualisation des KPIs et graphiques.
    \item \textbf{UC05 -- Import :} Traitement fichiers CSV bancaires.
    \item \textbf{UC06 -- Export :} Portabilité des données (RGPD).
\end{itemize}

\begin{figure}[H]
    \centering
    \safefig{chapters/assets/images/usage_global.png}{0.92\textwidth}{usage_global.png}
    \caption{Diagramme de cas d'usage global — My Smart Budget}
    \label{fig:usage_global}
\end{figure}

\begin{jury}
\textbf{Points de vigilance contrôlés :}
\begin{itemize}
    \item Les cas d'usage critiques sont tous traçables vers des User Stories.
    \item Les flux alternatifs (erreurs) sont documentés dans les critères d'acceptation.
    \item La distinction des rôles (User vs Admin) est claire.
\end{itemize}
\end{jury}

% ============================================================
% DIAGRAMME DE CLASSES
% ============================================================

\section{Modélisation statique : Diagramme de Classes}

Ce diagramme structure les données manipulées par l'application et sert de référence pour le développement Back-end (Express.js) et la base de données.

\begin{figure}[H]
    \centering
    \safefig{chapters/assets/images/class-diagram-app-budget-v3.png}{0.95\textwidth}{class-diagram-app-budget-v3.png}
    \caption{Diagramme de classes UML détaillé — My Smart Budget}
    \label{fig:class-diagram}
\end{figure}

\begin{tcolorbox}[colback=blue!5!white,colframe=blue!60!black,title=\textbf{Mapping UML $\leftrightarrow$ Code (pg-promise)}]
Chaque classe UML est traduite en table PostgreSQL et requêtes pg-promise. Cela m'a évité des incohérences entre la théorie (UML) et ce que j'ai réellement codé :

\begin{lstlisting}[language=SQL]
-- Classe UML User -> Table PostgreSQL
CREATE TABLE users (
  id           SERIAL PRIMARY KEY,
  email        VARCHAR(255) UNIQUE NOT NULL,
  password_hash TEXT NOT NULL,
  role         VARCHAR(20) DEFAULT 'USER',
  currency     VARCHAR(10) DEFAULT 'EUR',
  created_at   TIMESTAMPTZ DEFAULT NOW()
);

-- Relations definies dans le diagramme UML
CREATE TABLE transactions (
  id           SERIAL PRIMARY KEY,
  user_id      INTEGER NOT NULL REFERENCES users(id) ON DELETE CASCADE,
  category_id  INTEGER NOT NULL REFERENCES categories(id)
);
\end{lstlisting}
\end{tcolorbox}

% ============================================================
% DIAGRAMMES DE SEQUENCE
% ============================================================

\section{Modélisation dynamique : Diagrammes de Séquence}

Les séquences détaillent les interactions temporelles entre le Frontend, l'API et la Base de données.

\subsection{Séquence 1 : Authentification (US-01)}

\begin{figure}[H]
    \centering
    \safefig{chapters/assets/images/seq-auth-login.png}{0.92\textwidth}{seq-auth-login.png}
    \caption{Diagramme de séquence — Flux de connexion utilisateur (Login)}
    \label{fig:seq-auth}
\end{figure}

\textbf{Déroulement nominal :}
\begin{enumerate}
    \item Saisie des identifiants sur le Frontend (Next.js).
    \item Appel API \texttt{POST /auth/login}.
    \item Le Backend vérifie le hash du mot de passe (bcrypt).
    \item Génération d'un token JWT (signé HS256).
    \item Retour \texttt{200 OK} et redirection vers le Dashboard.
\end{enumerate}

Pour l'utilisateur, cela se traduit par une connexion fluide en moins de 2 secondes, sans aucune manipulation technique de sa part.

\subsection{Séquence 2 : Création de Transaction (US-03)}

\begin{figure}[H]
    \centering
    \safefig{chapters/assets/images/seq-create-transaction.png}{0.92\textwidth}{seq-create-transaction.png}
    \caption{Diagramme de séquence — Flux de création d'une transaction}
    \label{fig:seq-transaction}
\end{figure}

Ce diagramme illustre la validation des données par le \texttt{TransactionController} Express avant l'insertion en base via pg-promise, garantissant l'intégrité des données financières.

Pour l'utilisateur, cela signifie qu'une transaction mal saisie est rejetée immédiatement, avec un message d'erreur clair, avant tout enregistrement.

\subsection{Séquence 3 : Alertes Budgétaires (US-07)}

\begin{figure}[H]
    \centering
    \safefig{chapters/assets/images/seq-budget-alert.png}{0.92\textwidth}{seq-budget-alert.png}
    \caption{Diagramme de séquence — Flux automatique de déclenchement d'alertes}
    \label{fig:seq-alert}
\end{figure}

% ============================================================
% UI / UX
% ============================================================

\section{Conception de l'Interface (UI/UX)}

La conception graphique de \textbf{My Smart Budget} a suivi une démarche rigoureuse en trois niveaux de fidélité progressive : \textbf{Zoning} (structure fonctionnelle), \textbf{Wireframes} (mise en page détaillée) et \textbf{Maquettes Figma} (rendu haute fidélité). Cette approche \textbf{"Mobile First"} garantit une expérience cohérente sur tous les supports.

% ------------------------------------------------------------
\subsection{Méthodologie UI/UX appliquée}
% ------------------------------------------------------------

\begin{center}
\small
\begin{tabular}{|p{0.12\textwidth}|p{0.35\textwidth}|p{0.42\textwidth}|}
\hline
\textbf{Étape} & \textbf{Livrable} & \textbf{Objectif} \\
\hline
\textbf{1 — Zoning} & Schéma de découpage en zones & Valider la hiérarchie de l'information et la disposition des blocs fonctionnels \\
\hline
\textbf{2 — Wireframe} & Maquette basse fidélité & Définir la structure précise, les interactions et les contenus sans couleurs ni styles \\
\hline
\textbf{3 — Maquette HF} & Figma haute fidélité & Intégrer la charte graphique, les composants du Design System et les états interactifs \\
\hline
\end{tabular}
\end{center}

% ------------------------------------------------------------
\subsection{Page d'inscription (Register)}
% ------------------------------------------------------------

\subsubsection{Zoning — Structure fonctionnelle}

Le zoning de la page d'inscription (identifiant \textbf{Z-01}) définit cinq zones fonctionnelles empilées verticalement, conformément à l'approche Mobile First :

\begin{figure}[H]
    \centering
    \safefig{chapters/assets/images/zoning-register.png}{0.55\textwidth}{zoning-register.png}
    \caption{Zoning Z-01 — Page d'inscription (Register)}
    \label{fig:zoning-register}
\end{figure}

\begin{center}
\small
\begin{tabular}{|p{0.22\textwidth}|p{0.65\textwidth}|}
\hline
\textbf{Zone} & \textbf{Contenu et rôle} \\
\hline
\textbf{Header} & Logo de l'application, lien de retour vers la page de connexion \\
\hline
\textbf{Titre} & Titre principal « Create Account » + sous-titre de réassurance \\
\hline
\textbf{Formulaire} & Champs : Nom complet, Adresse e-mail, Mot de passe, Confirmation + case CGU \\
\hline
\textbf{CTA Create Account} & Bouton d'action principal — soumission du formulaire \\
\hline
\textbf{CTA Google / Apple} & Connexion sociale — alternative rapide sans saisie de mot de passe \\
\hline
\end{tabular}
\end{center}

\subsubsection{Wireframe haute fidélité}

\begin{figure}[H]
    \centering
    \safefig{chapters/assets/images/wireframe-register.png}{0.62\textwidth}{wireframe-register.png}
    \caption{Maquette haute fidélité — Page d'inscription}
    \label{fig:wireframe-register}
\end{figure}

La maquette finale reprend fidèlement le zoning et y ajoute :
\begin{itemize}
    \item Un formulaire en deux colonnes pour les champs \textit{Password} et \textit{Confirm}, optimisant l'espace sur écran.
    \item Un bouton « Create Account » à fond sombre (\texttt{\#2D3748}) pour un fort contraste visuel.
    \item Une section « Or Continue With » proposant deux alternatives d'authentification sociale (Google, Apple).
    \item Un footer minimaliste avec le copyright \textit{MSB The Ledger}.
\end{itemize}

\begin{jury}
\textbf{Décision de conception :} L'obligation de cocher les CGU avant de soumettre le formulaire répond à une exigence RGPD — l'utilisateur donne un consentement explicite avant la création de son compte.
\end{jury}

% ------------------------------------------------------------
\subsection{Structure générale des pages — Systèmes de grille}
% ------------------------------------------------------------

\subsubsection{Grille mobile (layout vertical)}

\begin{figure}[H]
    \centering
    \safefig{chapters/assets/images/wireframe-mobile-layout.png}{0.35\textwidth}{wireframe-mobile-layout.png}
    \caption{Schéma de grille mobile — Structure des pages principales}
    \label{fig:wireframe-mobile-layout}
\end{figure}

Sur mobile, toutes les pages de l'application partagent la même structure en colonne unique :
\begin{itemize}
    \item \textbf{Header} : Logo MSB + icônes de notifications et de profil utilisateur.
    \item \textbf{Titre} : Titre de la page courante.
    \item \textbf{Formulaire / Main Content Canvas} : Zone de contenu principale (graphiques, listes, formulaires).
    \item \textbf{Side Module A \& B} : Modules latéraux reformatés en blocs empilés sur mobile.
    \item \textbf{CTA Primary / Secondary} : Boutons d'action accessibles en bas de page.
    \item \textbf{Liens de navigation} : Grid 2×2 de liens secondaires.
\end{itemize}

\subsubsection{Grille desktop (layout multi-colonnes)}

\begin{figure}[H]
    \centering
    \safefig{chapters/assets/images/zoning-desktop.png}{0.92\textwidth}{zoning-desktop.png}
    \caption{Zoning desktop — Structure multi-colonnes des pages principales}
    \label{fig:zoning-desktop}
\end{figure}

Sur desktop, la grille adopte une structure à trois colonnes :
\begin{itemize}
    \item \textbf{Colonne gauche (Z-03)} : Barre de navigation latérale fixe avec les liens principaux (Overview, Transactions, Budgets, Savings).
    \item \textbf{Zone centrale (Z-02)} : Contenu principal — graphiques, tableaux, formulaires.
    \item \textbf{Colonne droite (Z-04)} : Modules contextuels — résumés, alertes, actions rapides.
    \item \textbf{Header (Z-01)} : Bandeau supérieur avec logo, navigation globale et profil.
\end{itemize}

Cette grille adaptative garantit que le contenu affiché sur mobile (colonnes empilées) et sur desktop (colonnes côte-à-côte) reste identique — seule la disposition spatiale change.

% ------------------------------------------------------------
\subsection{Page Dashboard}
% ------------------------------------------------------------

\subsubsection{Dashboard Desktop}

\begin{figure}[H]
    \centering
    \safefig{chapters/assets/images/wireframe-dashboard-desktop.png}{0.92\textwidth}{wireframe-dashboard-desktop.png}
    \caption{Wireframe — Dashboard Desktop (Smart Budget)}
    \label{fig:wireframe-dashboard-desktop}
\end{figure}

Le tableau de bord desktop est structuré en deux zones principales :
\begin{itemize}
    \item \textbf{Navigation latérale} : Menu fixe avec quatre entrées — \textit{Overview}, \textit{Transactions}, \textit{Budgets}, \textit{Savings} — et un accès aux \textit{Settings} en bas.
    \item \textbf{Zone centrale} : Deux rangées de modules graphiques (graphiques de synthèse + modules contextuels), suivies du tableau \textit{Recent Ledger} affichant les dernières transactions avec entité, catégorie, date et montant.
\end{itemize}

\subsubsection{Zoning — Dashboard Mobile}

\begin{figure}[H]
    \centering
    \safefig{chapters/assets/images/zoning-dashboard-mobile.png}{0.38\textwidth}{zoning-dashboard-mobile.png}
    \caption{Zoning — Dashboard Mobile (structure des zones)}
    \label{fig:zoning-dashboard-mobile}
\end{figure}

Le zoning du dashboard mobile décompose la page en zones empilées selon leur priorité d'affichage :

\begin{center}
\small
\begin{tabular}{|p{0.18\textwidth}|p{0.69\textwidth}|}
\hline
\textbf{Zone} & \textbf{Contenu et rôle} \\
\hline
\textbf{Z-01} & Header — Logo MSB + icône profil utilisateur \\
\hline
\textbf{Z-02 (large)} & KPI principal — Solde total (Total Balance) \\
\hline
\textbf{Z-02 × 2} & Modules secondaires côte-à-côte — Revenus et Dépenses du mois \\
\hline
\textbf{Z-02 (pleine largeur)} & Graphique de synthèse mensuelle \\
\hline
\textbf{Z-02 × 2} & Category Breakdown — deux catégories principales \\
\hline
\textbf{Z-02 (pleine largeur)} & Troisième catégorie (ligne complète) \\
\hline
\textbf{Z-02 (liste)} & Recent Transactions — liste des dernières opérations \\
\hline
\textbf{Z-02} & Bouton « View All » — lien vers la page Transactions \\
\hline
\textbf{Z-05} & Footer de navigation \\
\hline
\end{tabular}
\end{center}

\subsubsection{Wireframe — Dashboard Mobile}

\begin{figure}[H]
    \centering
    \safefig{chapters/assets/images/wireframe-dashboard-mobile.png}{0.42\textwidth}{wireframe-dashboard-mobile.png}
    \caption{Wireframe — Dashboard Mobile (MSB)}
    \label{fig:wireframe-dashboard-mobile}
\end{figure}

Sur mobile, le dashboard condense l'information en une colonne scrollable :
\begin{itemize}
    \item \textbf{Carte Total Balance} : Solde total affiché en grand, avec \textit{Monthly Income} et \textit{Monthly Expenses} en sous-titres.
    \item \textbf{Category Breakdown} : Grille 2×1 + ligne complète pour les catégories prioritaires (Shopping, Food \& Drink, Transport) avec leur pourcentage consommé.
    \item \textbf{Recent Transactions} : Liste des trois dernières opérations avec icône, nom du vendeur, horodatage et montant coloré (vert pour les crédits, rouge pour les débits).
\end{itemize}

% ------------------------------------------------------------
\subsection{Wireframes existants — Transactions et Ajout}
% ------------------------------------------------------------

\begin{figure}[H]
    \centering
    \begin{minipage}{0.48\textwidth}
        \centering
        \safefig{chapters/assets/images/wireframe-dashboard.png}{\textwidth}{wireframe-dashboard.png}
        \caption{Wireframe basse fidélité — Dashboard (version initiale)}
        \label{fig:wireframe-dashboard}
    \end{minipage}\hfill
    \begin{minipage}{0.48\textwidth}
        \centering
        \safefig{chapters/assets/images/wireframe-transactions.png}{\textwidth}{wireframe-transactions.png}
        \caption{Wireframe — Page Transactions}
        \label{fig:wireframe-transactions}
    \end{minipage}
\end{figure}

Ces wireframes basse fidélité ont constitué la première itération de conception, validant les flux utilisateurs avant la production des maquettes haute fidélité dans Figma.

\subsection{Charte Graphique (Design System)}

\begin{center}
\small
\begin{tabular}{|p{0.25\textwidth}|p{0.2\textwidth}|p{0.45\textwidth}|}
\hline
\textbf{Élément} & \textbf{Valeur} & \textbf{Usage} \\
\hline
\textbf{Bleu Primaire} & \texttt{\#1D4ED8} & Actions principales, Header \\
\hline
\textbf{Vert Succès} & \texttt{\#22C55E} & Revenus, Soldes positifs \\
\hline
\textbf{Orange Alerte} & \texttt{\#F97316} & Warning (Budget > 80\%) \\
\hline
\textbf{Rouge Danger} & \texttt{\#EF4444} & Erreurs, Dépenses, Dépassements \\
\hline
\textbf{Typographie} & \textit{Inter} & Police sans-serif moderne et lisible \\
\hline
\end{tabular}
\end{center}

% ============================================================
% BASE DE DONNÉES
% ============================================================

\section{Conception de la Base de Données (PostgreSQL)}

Le modèle de données a été conçu selon la méthode Merise (MCD/MLD/MPD) puis implémenté physiquement sous PostgreSQL 15.

\subsection{Modèle Conceptuel (MCD)}

\begin{figure}[H]
    \centering
    \safefig{chapters/assets/images/mcd-app-budget-v3.png}{0.95\textwidth}{mcd-app-budget-v3.png}
    \caption{Modèle Conceptuel de Données (MCD) — entités et relations}
    \label{fig:mcd}
\end{figure}

\subsection{Implémentation Physique (SQL)}

Extrait du script de création des tables (écrit et exécuté via pg-promise au démarrage du serveur) :

\begin{lstlisting}[language=SQL]
-- Table des transactions avec contraintes
CREATE TABLE transactions (
  id SERIAL PRIMARY KEY,
  user_id INTEGER NOT NULL REFERENCES users(id) ON DELETE CASCADE,
  category_id INTEGER NOT NULL REFERENCES categories(id),
  date DATE NOT NULL,
  amount NUMERIC(12,2) NOT NULL,
  type TEXT NOT NULL CHECK (type IN ('DEBIT','CREDIT')),
  label TEXT NOT NULL,
  created_at TIMESTAMPTZ DEFAULT NOW()
);

-- Index de performance pour le Dashboard
CREATE INDEX idx_transactions_user_date ON transactions(user_id, date DESC);
\end{lstlisting}

% ============================================================
% ARCHITECTURE
% ============================================================

\section{Architecture Technique 3-Tiers}

L'application repose sur une séparation stricte des responsabilités.

\begin{figure}[H]
    \centering
    \safefig{chapters/assets/images/arch-3tiers-app-budget-v3.png}{0.92\textwidth}{arch-3tiers-app-budget-v3.png}
    \caption{Architecture 3-tiers : Next.js (Front) — Express (Back) — PostgreSQL (Data)}
    \label{fig:arch-3tiers}
\end{figure}

\begin{itemize}
    \item \textbf{Tier 1 (Présentation) :} \textbf{Next.js 15}. Rendu hybride (SSR/CSR), composants React, validation formulaires.
    \item \textbf{Tier 2 (Métier) :} \textbf{Node.js/Express}. API RESTful, middlewares JWT, logique métier, validation des requêtes.
    \item \textbf{Tier 3 (Données) :} \textbf{PostgreSQL 16} (pg-promise) + \textbf{MongoDB} (Mongoose). Données relationnelles et non-structurées, intégrité référentielle et schémas flexibles.
\end{itemize}

\section{Matrice de Cohérence Globale}

Le tableau ci-dessous démontre la traçabilité complète du besoin jusqu'au code :

\begin{center}
\small
\begin{tabular}{|p{0.15\textwidth}|p{0.15\textwidth}|p{0.2\textwidth}|p{0.2\textwidth}|p{0.15\textwidth}|}
\hline
\textbf{Use Case} & \textbf{Classe UML} & \textbf{Endpoint API} & \textbf{Table SQL} & \textbf{Écran UI} \\
\hline
UC01 Auth & User & \texttt{POST /auth/login} & \texttt{users} & Login \\
\hline
UC02 Transac. & Transaction & \texttt{POST /transactions} & \texttt{transactions} & Formulaire \\
\hline
UC03 Budget & BudgetEnv. & \texttt{GET /budgets} & \texttt{budget\_envelopes} & Liste env. \\
\hline
UC04 Dash. & (Agregats) & \texttt{GET /analytics} & (Join queries) & Dashboard \\
\hline
UC05 Import & ImportJob & \texttt{POST /import} & \texttt{import\_jobs} & Upload \\
\hline
\end{tabular}
\end{center}

\begin{tcolorbox}[colback=green!5!white,colframe=green!60!black,title=\textbf{Exemple de traçabilité : UC02 Transactions}]
\textbf{Besoin :} US-03 "Créer une transaction" \\
$\rightarrow$ \textbf{UML :} Classe \texttt{Transaction}, Séquence \texttt{CreateTransaction} \\
$\rightarrow$ \textbf{API :} \texttt{TransactionsController.create()} \\
$\rightarrow$ \textbf{DB :} Table \texttt{transactions} (FK user\_id, category\_id) \\
$\rightarrow$ \textbf{UI :} Page \texttt{/transactions/new}
\end{tcolorbox}

\section{Ressources et Références}

\begin{itemize}
    \item \textbf{UML \& Merise :} \url{https://www.uml-diagrams.org/}
    \item \textbf{Frameworks :} \url{https://nextjs.org/} (Front), \url{https://expressjs.com/} (Back)
    \item \textbf{Base de données :} \url{https://www.postgresql.org/docs/}
    \item \textbf{Outils :} Figma (UI), Draw.io (Diagrammes)
\end{itemize}